\section{Задачи, выполнимые за $O(1)$}

\begin{frame}{Некоторые договоренности}
	Для удобства мы будем считать, что точки {\itshape в общем положении:}
	\begin{itemize}
		\item никакие 3 из них не лежат на одной прямой,
		\item никакие 4 из них не лежат на одной окружности,
		\item у них нет общих $x$ и $y$ координат.
	\end{itemize}
	В совокупности это верно {\itshape почти всегда.}
\end{frame}


\subsection{Применения аналитической геометрии}

\begin{frame}{Задание фигур уравнениями}
  При работе с геом. объектами удобно задавать их уравнениями
  \begin{itemize}
    \item Прямая:
      \[\ell\colon ax + by + c = 0 \text{\ \ или\ \ } \ell\colon y = kx + b.\]

    \item Прямая через точки $A(x_1, y_1)$, $B(x_2, y_2)$:
      \[a = y_2 - y_1,\ \ b = x_1 - x_2,\ \ c = -(ax_1 + by_1).\]

    \item Окружность:
      \[(x - x_0)^2 + (y - y_0)^2 = R^2 \text{:\ \ }
        (x_0, y_0) \text{ — центр,\ \ }
        R \text{ — радиус.}\]
  \end{itemize}
\end{frame}

\begin{frame}{Расстояние от точки до прямой}

	Если прямая задана как  $\ell\colon ax + by + c = 0$, то расстояние от точки $M(x_0, y_0)$ до неё можно рассчитать как
	\[ d(M, \ell) = \frac{\left| ax_0 + by_0 + c \right|}{\sqrt{a^2 + b^2}} \]

	\begin{center}
	\begin{tikzpicture}[scale = 0.58]
	\coordinate[label=left:$M$] (M) at (3,5);
	\draw[fill4, fill = fill4] (3,5) circle (0.07);
	\draw[fill4, fill = fill4] (4,1) circle (0.07);
	\draw[fill4, thick] (0, 0)--(8, 2) node[white, left, right = 0.05]{$\ell$};
	\draw[fill4, thick] (3, 5)--(4, 1);
	\end{tikzpicture}
	\end{center}

\end{frame}

\subsection{Центр описанной окружности}

\begin{frame}{Центр описанной окружности}

	\vspace{2.5mm}
	\begin{task}

		Дана тройка точек $a, b, c$. Требуется найти центр описанной окружности $\triangle abc$.

	\end{task}

	\begin{center}
		\begin{tikzpicture}[sty/.style={fill = fill4, circle, inner sep=2pt}]
		\coordinate[sty, label=left:$a$] (A) at (0,0); \coordinate[sty,label=right:$b$](B) at (3,1); \coordinate[sty,label=above:$c$](C) at (1,3);
		\draw[thick, fill4] (A)--(B)--(C)--(A);
		\coordinate (T) at ($(A)!1!60:(B)$); \draw [name path=h1, dashed]
		(T)--($(A)!(T)!(B)$); \coordinate (T) at ($(B)!1!60:(C)$);
		\draw [name path=h2, dashed] (T)--($(B)!(T)!(C)$); \draw [name intersections=
		{of=h1 and h2, by=O}]; \node[sty, label=right:$O$] at (O) {};
		\draw[thick, fill5] (O) let \p1=($(O)-(A)$)
		in circle ({veclen(\x1,\y1)});
	   \end{tikzpicture}
	\end{center}

\end{frame}


\begin{frame}{Центр описанной окружности}
  \begin{enumerate}
    \item Серединные перпендикуляры к $[ac]$ и $[bc]$: \begin{itemize}
      \item произведение угловых коэф-тов перп. прямых — \(-1\),
      \item середина отрезка — полусумма координат.
    \end{itemize} \medskip
    
    \item Пересечение прямых~— решение линейной системы.
  \end{enumerate} \bigskip

  Другие замечательные точки треугольника находятся аналогично.
\end{frame}


\subsection{Ориентация и косое произведение}

\begin{frame}{Угол между векторами}

	Рассмотрим вектора $\vec{a} = (x_1, y_1)$ и $\vec{b} = (x_2, y_2)$.
	\[ \langle \vec{a}, \vec{b} \rangle = x_1 x_2 + y_1 y_2 = |\vec{a}||\vec{b}| \cos{\angle{(\vec{a}, \vec{b})}} = \sqrt{x_1^2 + y_1^2}\sqrt{x_2^2 + y_2^2} \cos{\angle(\vec{a}, \vec{b})}.\]
	Значит, мы знаем, как найти угол между векторами $\vec{a}$ и $\vec{b}$:
	\[ \angle(\vec{a}, \vec{b}) = \arccos\lr*{\frac{x_1 y_1 + x_2 y_2}{\sqrt{x_1^2 + y_1^2} \sqrt{x_2^2 + y_2^2}}} \]

\end{frame}

\begin{frame}{Косое произведение}
  \begin{defn}
    {\itshape Косое произведение} векторов $\vec{a} = (x_1, y_1)$ и $\vec{b} = (x_2, y_2)$:
      \[ \vec{a} \wedge \vec{b} =  \begin{vmatrix} x_1 & y_1 \\ x_2 & y_2 \end{vmatrix}
        = x_1 y_2 - x_2 y_1.\]
  \end{defn}

  \begin{thm}
    \[\vec{a} \wedge \vec{b} = |\vec{a}| |\vec{b}| \sin \lr*{\angle \lr*{\vec{a},\vec{b}}},\]
    $\angle\lr*{\vec{a}, \vec{b}}$~—
    угол вращения против часовой стрелки от $\vec{a}$ к $\vec{b}$.
  \end{thm}  \vspace{-1cm}
\end{frame}


\begin{frame}{Площадь треугольника}
\begin{thm}
  $\vec{a} \wedge \vec{b} = 2 S_{\triangle{ABD}}$, где площадь {\itshape ориентированная.}

	\begin{center} \begin{tikzpicture}
	\coordinate[label=left:$A$] (A) at (0,0); \coordinate[label=left:$B$] (B) at (1,3); \coordinate[label=right:$C$] (C) at (4,3); \coordinate[label=right:$D$] (D) at (3,0);
	\draw[->, fill4, thick] (0,0)--(1,3) node[midway, left = 0.1]{$\vec{a}$};
	\draw[->, fill4, thick] (0,0)--(3,0) node[midway, below = 0.1]{$\vec{b}$};
	\draw[fill4, thick] (1,3)--(3,0);
	\draw[fill4, fill=fill4, fill opacity=0.5, draw opacity = 0] plot coordinates{(0,0) (1,3) (3, 0)}--cycle;
	\draw (A)--(B)--(C)--(D)--cycle;
	\end{tikzpicture} \end{center}
\end{thm} \vspace{-0.8cm}
\end{frame}


\begin{frame}{Ориентация}

	\begin{defn}

		Будем говорить, что тройка точек $(a, b, c)$ {\itshape положительно ориентирована} и писать $\sign(a, b, c) > 0$, если поворот вектора
		$\overrightarrow{ba}$ к вектору $\overrightarrow{bc}$ осуществляется против часовой стрелки.

	\end{defn}

	\begin{block}{Замечение}

		Ориентация тройки точек $(a, b, c)$ совпадает со знаком косого произведения $\overrightarrow{ba} \wedge \overrightarrow{bc}$.

	\end{block}
\end{frame}

\begin{frame}{Ориентация}

	\begin{center}
	\begin{tikzpicture}
	\begin{scope}
	\coordinate[label= left:$b$] (b) at (0,1); \coordinate[label= left :$c$] (c) at (1,4); \coordinate[label=right:$a$] (a) at (3,1);
	\draw[->, fill4, thick] (0,1)--(1,4) node[midway, left = 0.1]{$\vec{b}$};
	\draw[->, fill4, thick] (0,1)--(3,1) node[midway, below = 0.1]{$\vec{a}$};
	\draw pic[draw, ->, thick, fill1, angle radius=1cm, shift={(6mm,1mm)}] {angle=a--b--c};
	\node[draw] at (1,5) {$\sign(a, b, c) > 0$};
	\end{scope}

	\begin{scope}[xshift=5cm]
	\coordinate[label= left:$b$] (b) at (0,1); \coordinate[label= left :$a$] (a) at (1,4); \coordinate[label=right:$c$] (c) at (3,1);
	\draw[->, fill4, thick] (0,1)--(1,4) node[midway, left = 0.1]{$\vec{a}$};
	\draw[->, fill4, thick] (0,1)--(3,1) node[midway, below = 0.1]{$\vec{b}$};
	\draw pic[draw, <-,  thick, fill1, angle radius=1cm, shift={(6mm,1mm)}] {angle=c--b--a};
	\node[draw] at (1,5) {$\sign(a, b, c) < 0$};
	\end{scope}
	\end{tikzpicture}
	\end{center}

\end{frame}

\subsection{Пересечение отрезков}

\begin{frame}{Пересечение отрезков}
  \begin{task}
    Дана четверка точек $a, b, c, d$. Требуется определить, \\
    пересекаются ли отрезки $[ab]$ и $[cd]$.
  \end{task}

\begin{tabular}{ll}
  \makecell[l]{
    Отрезки $[ab]$ и $[cd]$ пересекаются \\ т. и т. т., когда: \\
    \ \ • концы $a, b$ лежат по разные стороны \\ \ \ от прямой $(cd)$, \\
    \ \ • концы $c, d$ лежат по разные стороны \\ \ \ от прямой $(ab)$. \\
    Это четыре проверки знаков.
  } &
  \makecell[l]{\hspace{1cm}
    \begin{tikzpicture}[scale=0.89]
      \coordinate[label=left:$c$] (c) at (0,0);
      \coordinate[label=left:$d$] (d) at (3,3);
      \coordinate[label=right:$a$] (a) at (0,4);
      \coordinate[label=right:$b$] (b) at (3,1);
      %
      \draw[thick,fill5] (b)--(c)--(a);
      \fill[fill4] (c) circle [radius=0.7mm]; \fill[fill4] (d) circle [radius=0.7mm];
      \fill[fill4] (a) circle [radius=0.7mm]; \fill[fill4] (b) circle [radius=0.7mm];
      \draw[thick,fill4] (a)--(b) (c)--(d);
      %
      \draw pic[draw,->,thick,fill1,angle radius=0.8cm] {angle=b--c--d};
      \draw pic[draw,<-,thick,fill1,angle radius=0.8cm] {angle=d--c--a};
    \end{tikzpicture}
  }
\end{tabular}  \vspace{-0.8cm}
\end{frame}


\begin{frame}{Пересечение отрезков}
	\begin{center}
	\begin{algorithmic}
	\State \underline{\textsc{Intersect}(a, b, c, d):}
	\If{$\sign(a, c, d) = \sign(b, c, d)$}
		\State \Return \textsc{False}
	\Else
    	\If{$\sign(a, b, c) = \sign(a, b, d)$}
        	\State \Return \textsc{False}
	\Else
		\State \Return \textsc{True}
    	\EndIf
	\EndIf
	\end{algorithmic}
	\end{center}  \vspace{-0.8cm}
\end{frame}
